% Options for packages loaded elsewhere
\PassOptionsToPackage{unicode}{hyperref}
\PassOptionsToPackage{hyphens}{url}
\documentclass[
  ignorenonframetext,
]{beamer}
\newif\ifbibliography
\usepackage{pgfpages}
\setbeamertemplate{caption}[numbered]
\setbeamertemplate{caption label separator}{: }
\setbeamercolor{caption name}{fg=normal text.fg}
\beamertemplatenavigationsymbolsempty
% remove section numbering
\setbeamertemplate{part page}{
  \centering
  \begin{beamercolorbox}[sep=16pt,center]{part title}
    \usebeamerfont{part title}\insertpart\par
  \end{beamercolorbox}
}
\setbeamertemplate{section page}{
  \centering
  \begin{beamercolorbox}[sep=12pt,center]{section title}
    \usebeamerfont{section title}\insertsection\par
  \end{beamercolorbox}
}
\setbeamertemplate{subsection page}{
  \centering
  \begin{beamercolorbox}[sep=8pt,center]{subsection title}
    \usebeamerfont{subsection title}\insertsubsection\par
  \end{beamercolorbox}
}
% Prevent slide breaks in the middle of a paragraph
\widowpenalties 1 10000
\raggedbottom
\AtBeginPart{
  \frame{\partpage}
}
\AtBeginSection{
  \ifbibliography
  \else
    \frame{\sectionpage}
  \fi
}
\AtBeginSubsection{
  \frame{\subsectionpage}
}
\usepackage{iftex}
\ifPDFTeX
  \usepackage[T1]{fontenc}
  \usepackage[utf8]{inputenc}
  \usepackage{textcomp} % provide euro and other symbols
\else % if luatex or xetex
  \usepackage{unicode-math} % this also loads fontspec
  \defaultfontfeatures{Scale=MatchLowercase}
  \defaultfontfeatures[\rmfamily]{Ligatures=TeX,Scale=1}
\fi
\usepackage{lmodern}
\ifPDFTeX\else
  % xetex/luatex font selection
\fi
% Use upquote if available, for straight quotes in verbatim environments
\IfFileExists{upquote.sty}{\usepackage{upquote}}{}
\IfFileExists{microtype.sty}{% use microtype if available
  \usepackage[]{microtype}
  \UseMicrotypeSet[protrusion]{basicmath} % disable protrusion for tt fonts
}{}
\makeatletter
\@ifundefined{KOMAClassName}{% if non-KOMA class
  \IfFileExists{parskip.sty}{%
    \usepackage{parskip}
  }{% else
    \setlength{\parindent}{0pt}
    \setlength{\parskip}{6pt plus 2pt minus 1pt}}
}{% if KOMA class
  \KOMAoptions{parskip=half}}
\makeatother
\usepackage{longtable,booktabs,array}
\newcounter{none} % for unnumbered tables
\usepackage{calc} % for calculating minipage widths
\usepackage{caption}
% Make caption package work with longtable
\makeatletter
\def\fnum@table{\tablename~\thetable}
\makeatother
\setlength{\emergencystretch}{3em} % prevent overfull lines
\providecommand{\tightlist}{%
  \setlength{\itemsep}{0pt}\setlength{\parskip}{0pt}}
\usepackage{bookmark}
\IfFileExists{xurl.sty}{\usepackage{xurl}}{} % add URL line breaks if available
\urlstyle{same}
\hypersetup{
  hidelinks,
  pdfcreator={LaTeX via pandoc}}

\author{\texorpdfstring{}{}}
\date{}

\begin{document}

\section{Transatlantic Architectures of Monetary Collapse: Britain and
America
Compared}\label{transatlantic-architectures-of-monetary-collapse-britain-and-america-compared}

\begin{frame}{The Divergent Ratios and Their Structural Consequences}
\protect\phantomsection\label{the-divergent-ratios-and-their-structural-consequences}
The British and American bimetallic experiments, though sharing the same
architecture---two metals bound by statutory fiat---diverged in ways
that illuminate both the generic logic of bimetallic instability and the
specific political economies encoded in each nation's monetary regime. A
systematic comparison reveals how Gresham's Law operated with identical
structural logic but divergent political consequences across the
Atlantic.

{\def\LTcaptype{none} % do not increment counter
\begin{longtable}[]{@{}
  >{\raggedright\arraybackslash}p{(\linewidth - 4\tabcolsep) * \real{0.2037}}
  >{\raggedright\arraybackslash}p{(\linewidth - 4\tabcolsep) * \real{0.3704}}
  >{\raggedright\arraybackslash}p{(\linewidth - 4\tabcolsep) * \real{0.4259}}@{}}
\toprule\noalign{}
\begin{minipage}[b]{\linewidth}\raggedright
Parameter
\end{minipage} & \begin{minipage}[b]{\linewidth}\raggedright
Britain (1717--1816)
\end{minipage} & \begin{minipage}[b]{\linewidth}\raggedright
United States (1792--1873)
\end{minipage} \\
\midrule\noalign{}
\endhead
Statutory Ratio & 1:15.21 (Newton, 1717) & 1:15 (Hamilton, 1792);
revised to 1:16 (1834) \\
Market Reality & \textasciitilde1:14.5 (European) &
\textasciitilde1:15.5 (European, post-1792) \\
Direction of Drainage & Silver exported to France/India & Gold exported
to Europe (pre-1834); Silver exported (post-1834) \\
Central Bank & Bank of England (1694) & First Bank of US (1791--1811);
Second Bank (1816--1836) \\
Suspension Event & Privy Council order, 26 Feb 1797; Royal Assent, 3 May
1797 & Suspension of specie convertibility, 1862--1879 \\
Demonetization & Coinage Act 1816 (56 Geo. III c.68) & Fourth Coinage
Act, 12 Feb 1873 \\
Popular Resistance & Minimal (no Free Silver equivalent) & Free Silver
movement (1873--1900) \\
Prophetic Critique & Blake's \emph{Jerusalem}, \emph{Milton} & Bryan's
``Cross of Gold'' (1896) \\
\bottomrule\noalign{}
\end{longtable}
}

The table reveals a structural insight: under Hamilton's 15:1 ratio,
which \emph{undervalued} gold relative to the European market of
approximately 15.5:1, the United States initially experienced the mirror
image of Britain's drainage. American gold was exported to Europe where
it commanded higher value, while silver accumulated
domestically---making silver, paradoxically, the \emph{de facto}
American standard through the 1820s {[}@laughlin1898bimetallism;
@officer2001bimetallism{]}. This is precisely the inverse of Britain's
experience under Newton's 15.21:1 ratio, which overvalued gold and
drained silver. The 1834 revision to approximately 16:1 reversed this
dynamic, deliberately overvaluing gold to attract it into domestic
circulation---but at the cost of triggering the same outward silver
flight that had devastated British coinage a century earlier
{[}@redish1990evolution{]}.
\end{frame}

\begin{frame}{Jackson's War on the Bank: The Limits of Institutional
Destruction}
\protect\phantomsection\label{jacksons-war-on-the-bank-the-limits-of-institutional-destruction}
Andrew Jackson's veto of the Second Bank of the United States in 1832,
followed by his executive removal of federal deposits in 1833 and the
Bank's ultimate demise in 1836, constitute an episode of resonance for
understanding the structural limits of monetary reform. Jackson, the
agrarian populist from the frontier, declared the Bank a ``monopoly''
and a ``hydra of corruption'' that enriched Eastern financial elites at
the expense of the producing classes {[}@howe2007what;
@wilentz2005rise{]}. Daniel Walker Howe's \emph{What Hath God Wrought}
demonstrates that Jackson's Bank War was not merely an economic dispute
but a fundamental contest over the meaning of democratic sovereignty in
the new republic: who should control the creation of money, and to whose
benefit {[}@howe2007what{]}?

In Blake's mythological terms, this is Orc's revolutionary assault on
the Urizenic institution---the fiery destruction of the centralized
architecture of credit. Yet Jackson's destruction of the Bank unleashed
a ``Free Banking'' era of wildcat state-chartered institutions that
issued unregulated paper currency---precisely the proliferating
``allegoric riches'' Blake treats in \hyperlink{sec-allegoric-riches}{§
The fixing of labour and allegoric riches}, now scaled horizontally
across state charters. Sean Wilentz documents how the period between the
Bank's destruction (1836) and the Civil War saw over 1,600
state-chartered banks issuing their own notes, producing a fragmented,
unreliable, and frequently fraudulent currency system
{[}@wilentz2005rise{]}.

The Jacksonian paradox recapitulates the Blakean diagnosis with
fidelity: destroying the centralized institution does not automatically
produce liberation. In the language of \emph{America a Prophecy}, where
Orc's revolutionary fire has consumed the old order, Urizen descends
again---reconstituting himself in new institutional forms. This cycle of
revolution and reconsolidation is not a failure of the revolutionary
impulse but a consequence of insufficient cognitive transformation---a
structural feature of Blake's prophetic epistemology
{[}@friedman2026blake; @friedman2026doors{]}. Without the integrative
Fourfold Vision that coordinates all four Zoas---reason, passion,
sensation, and imagination---the destruction of one Urizenic institution
merely clears the ground for another. As Blake writes in \emph{America a
Prophecy}:

\begin{quote}
``Urizen, who sat / Above all heavens, in thunders wrapp'd, emerg'd his
leprous head From out his holy shrine, his tears in deluge piteous
Falling into the deep sublime'' (\emph{America a Prophecy}, Plates
16--17; Erdman pp.~57--58) {[}@erdman1988complete; @blake1793america{]}
\end{quote}

Revolution alone is insufficient; the Urizenic principle merely
reconstitutes itself in new institutional forms. Jackson dismantled the
\emph{federal} monopoly on credit only to democratize the issuance of
unbacked paper through state banks---a horizontal proliferation of the
same abstracting technology. Without the qualitative vision that
transcends both metallic fetishism and paper abstraction, the
destruction of one tyranny merely inaugurates another.
\end{frame}

\begin{frame}{The Greenback Suspension: America's Unprecedented Wartime
Monetary Swerve}
\protect\phantomsection\label{the-greenback-suspension-americas-unprecedented-wartime-monetary-swerve}
The American Civil War (1861--1865) produced its own swerve into fiat
abstraction, paralleling the British Bank Restriction of 1797. In 1862,
the United States Treasury began issuing Legal Tender Notes
(``Greenbacks'')---unbacked paper currency designed to finance the Union
war effort. Like the Bank Restriction, the American suspension of specie
convertibility severed the connection between the circulating medium and
any metallic anchor, plunging the economy into a void of fiat
abstraction that persisted until the Specie Resumption Act of 1875
restored gold convertibility on 1 January 1879
{[}@laughlin1898bimetallism; @unger1964greenback{]}.

Irwin Unger's Pulitzer Prize-winning \emph{The Greenback Era}
demonstrates that the suspension was far more than a wartime expedient:
it catalyzed a reorganization of American political coalitions around
the question of what money \emph{is} {[}@unger1964greenback{]}. The
``soft money'' faction---farmers, debtors, Western
entrepreneurs---demanded the permanent retention of Greenbacks as an
inflationary counterweight to the deflationary gold standard. The ``hard
money'' faction---Eastern bankers, creditors, industrial
capitalists---insisted on resumption and gold convertibility. This
division prefigured and structurally anticipated the Free Silver
movement of the 1880s--1890s.

The structural parallel with Britain is precise: both the 1797 British
suspension and the 1862 American suspension were triggered by
existential military crises; both produced inflationary spirals that
devastated the laboring classes while enriching financial speculators;
and both were resolved by a return to metallic standards that imposed
deflationary austerity upon the producing classes. In both cases,
Blake's prophetic analysis holds: the oscillation between the crushing
geometry of metallic standards and fiat paper constitutes a double
bind---a pendulum swing between two modes of quantification, neither of
which addresses the fundamental crisis of human value.
\end{frame}

\begin{frame}{From Albion to America: Prophetic and Monetary
Transmission}
\protect\phantomsection\label{from-albion-to-america-prophetic-and-monetary-transmission}
Blake's \emph{America a Prophecy} concludes not with Orc's triumph but
with Urizen's renewed assault---his ``icy magazine'' poured upon the
Atlantic to contain the revolutionary fire for ``twelve years'' before
``France receiv'd the Demon's light'' (Plates 16--17; Erdman pp.~57--58)
{[}@erdman1988complete; @blake1793america{]}. This prophetic
temporality---revolution contained but not extinguished, revolutionary
energy transmitted across borders---maps with precision onto the
transmission of bimetallic crises across the Atlantic. Britain's 1816
demonetization of silver preceded, and structurally anticipated,
America's 1873 demonetization. Barry Eichengreen's \emph{Globalizing
Capital} traces this convergence as part of a broader ``network effect''
in which the gold standard spread not through inherent superiority but
through the gravitational pull of London's capital markets, which made
gold convertibility the price of access to international credit
{[}@eichengreen2019globalizing{]}.

This transatlantic transmission functions as a key instance of what
Peirce calls the ``community of inquirers''---the recognition that no
single nation's monetary experiment operates in isolation, but that all
are bound in a web of mutual inference and structural feedback
{[}@friedman2026blake{]}. Blake's prophetic vision anticipated this
insight: Orc's fire cannot be confined to a single continent, and
Urizen's counterrevolution is likewise global. The arc from Newton's
1717 ratio through Hamilton's 1792 ratio, through both nations'
suspensions and ultimate demonetizations of silver, traces a single
structural trajectory: the global establishment of monometallic
quantification over bimetallic generation. Yet the 19th-century gold
standard was merely a transitional phase toward a more radical monetary
abstraction.
\end{frame}

\end{document}
